\section{The coefficient-trace comparison}\label{sec:c-comparison}

This section records the C argument for the Johnson replacement collar.

\subsection{The product Hattori equivalence}

The compact words give the letter counts
\[
p_y=42+2=44,\qquad p_z=269+2=271.
\]
In the cyclic $m_2$ word, every one of the 42
$y/Y$ events is followed immediately by a distinct $z/Z$ event; the same
holds for the two $y/Y$ events of $r_{xy}$.  The witness lists all 44 pairs.
Those pairs are realized as 44 product ribbons of the P0 control
strands along their recorded normals, with $227$ model $z$-circles placed
off the P0 cube.  The scope is the Johnson replacement collar, not an
embedded cut in $\partial W_2$.

Let $\mathcal C$ be the pregraded rational tangle category with objects
$T:P_{86}\to P_{88}$, and let $F$ be the framed west-to-east tangle formed
by those rectangles.  MWW's one-handle theorem identifies the cut coefficient
as
\[
M_R(T,T')=\KhR_2(R\cup T'\cup\overline T)
\]
with its displayed shift and gluing actions
\cite[Theorem~4.7]{ManolescuWalkerWedrich2023}.  Pivotal tangle duality and
K\"unneth over $\Q$ give graded isomorphisms on this replacement collar
\begin{equation}\label{eq:HattoriActual}
H_{T,T'}:M_R(T,T')\xrightarrow{\cong}
\operatorname{Hom}_{\mathcal C}(FT,FT')\{-44\}
\otimes A^{\otimes227},
\qquad A=\Q[X]/(X^2).
\end{equation}
These maps form a coefficient-bimodule isomorphism on the replacement collar.
Gluing $f:S\to T$ before the source link and then applying the product
isotopy is isotopic rel boundary to first applying the isotopy and
precomposing by $Ff$.  The right action is the same argument with
postcomposition.  Thus both action squares are associativity of gluing for
one simultaneous surface.  At the chain level the two squares are
\[
\begin{array}{ccc}
C_R(T,T')&\xrightarrow{\ H_{T,T'}\ }&
C(FT,FT')\otimes C(S^1)^{\otimes227}\\
\big\downarrow {L_f}&&\big\downarrow {L_{Ff}\otimes1}\\
C_R(S,T')&\xrightarrow{\ H_{S,T'}\ }&
C(FS,FT')\otimes C(S^1)^{\otimes227},
\end{array}
\qquad
\begin{array}{ccc}
C_R(T,T')&\xrightarrow{\ H_{T,T'}\ }&
C(FT,FT')\otimes C(S^1)^{\otimes227}\\
\big\downarrow {R_g}&&\big\downarrow {R_{Fg}\otimes1}\\
C_R(T,U)&\xrightarrow{\ H_{T,U}\ }&
C(FT,FU)\otimes C(S^1)^{\otimes227}.
\end{array}
\tag{17}
\]
Here $C_R$ and $C$ are the strict Blanchet--Khovanov foam complexes.  Both
composites in each square are the same glued foam after exchanging two
disjoint collars.  BHPW strict functoriality removes the projective sign, and
taking homology gives \eqref{eq:HattoriActual} and both MWW action squares
for the Johnson replacement collar.  Equation~\eqref{eq:HattoriActual} is
not a chain-level complex of an embedded cut in $\partial W_2$.

The ordinary quotient statement, including both cyclic-relation spans, is
checked in the formal development.\footnote{See
\texttt{Smooth4PC/RepresentableCoefficient.lean}.}

\begin{lemma}[Actual coefficient bimodule]\label{lem:C1}
For the P0 replacement, \eqref{eq:HattoriActual} is a natural isomorphism of
the actual MWW coefficient bimodule, with both gluing actions.  Moreover the
map \(\operatorname{Sh}_q\) below is the quantum trace of that bimodule, not a
dimension-matched substitute.
\end{lemma}

\begin{remark}[Status of Lemma~\ref{lem:C1}]
\Open.  Established: the $44$ model rectangles from the P0 control strands
with product-normal $z$-sides, giving the Johnson $y$-then-$z$ pairing and
$227$ model $z$-circles, and the representable $HH_0$ reduction of
Section~\ref{sec:exact-algebra}.  Not established: the isotopy
\[
R\cup\overline T\cup T'\;\cong\;\overline{FT}\circ FT'\sqcup 227\text{ circles}
\]
for the actual cut tangle $R$, fixing the boundary, the detector ball and
the gluing supports, and hence the chain-level action squares (17).  Without
that isotopy, \eqref{eq:HattoriActual} is a model isomorphism and the map
$\operatorname{Sh}_q$ below is not yet identified as the quantum trace of
the actual coefficient bimodule.
\end{remark}

\subsection{Quantum trace and the endpoint map}

Set $R_q=\Q[q,q^{-1}]$.  The pregraded cyclic relation is
\[
L_f(m)=q^{|f|}R_f(m).
\]
By Lemma~\ref{lem:C1}, \eqref{eq:HattoriActual} and the counits are homogeneous
coefficient morphisms on the Johnson replacement collar.  BPW's canonical vertical-to-horizontal functor and the BHPW strict
tangle/foam functor give
\begin{equation}\label{eq:actualShadow}
\operatorname{Sh}_q:q\operatorname{Tr}(\mathcal C;M_R)
\longrightarrow\operatorname{Hom}_{R_q}(E_{86},E_{88}),
\end{equation}
after the flat-base qHH$_0$/Chern identification
\cite{BeliakovaPutyraWehrli2019,BeliakovaHogancampPutyraWehrli2023}.  Here
\[
E_{86}=(V^{\otimes86})_{\lambda=86},\qquad
E_{88}=(V^{\otimes88})_{\lambda=86};
\]
the first has rank one and the second rank 88.  The tangle maps are
$U_q(\mathfrak{sl}_2)$ intertwiners before restriction to these weight
spaces.

Base change $q\mapsto1+h$ factors through rational localization and the
completion of a Noetherian local ring, hence is flat.  Reduction modulo $h$
of the trace cokernel simply replaces $q^{|f|}$ by one, giving the ordinary
MWW coefficient $HH_0$; right exactness of tensor product suffices for this
last assertion.

\subsection{Selected class and divided detector}

Let $U:P_{86}\to P_{88}$ be the selected one-cup tangle in the P0 endpoint
order and put $T_1=F^{-1}U$.  Define
\[
v_T=H_{T_1,T_1}^{-1}
  (\operatorname{Id}_U\otimes X^{\otimes227}).
\]
The circle counits give
$\operatorname{Sh}_h([v_T])=u_h$.  Its absolute degree is
\[
-44+227+315-4=494.
\]

The normalized checked R-matrix on adjacent one-defect basis vectors is
\[
\begin{pmatrix}1-q^{-2}&q^{-1}\\q^{-1}&0\end{pmatrix}.
\]
After the position basis change and $t=q^{-2}$, it is the public positive
Burau block; the negative block is its inverse.  The physical cable has mixed
orientations.  BPW's oriented-cabling formula identifies it with the
all-positive model up to a writhe-dependent grading shift and pivotal basis
maps.  The complete public word has writhe zero, so the global shift
vanishes.  Any degree-zero permutation or pivotal chart $P$ transports the
operator, cup and cap simultaneously:
\[
W\mapsto PWP^{-1},\qquad u\mapsto Pu,\qquad
\ell\mapsto\ell P^{-1}.
\]
The matrix coefficient is unchanged; we choose the public endpoint labels
after this common transport.  The constant terms are fixed by the single
collar endpoint convention of Section~\ref{sec:exact-algebra}: the selected
cup joins the physical endpoints bound to collar positions $2$ and $87$, and
BPW's cup/cap formulas give
\[
u_0=e_2-e_{87},\qquad
\ell_0=e_{87}^*-e_2^*,
\]
as in \eqref{eq:u} and \eqref{eq:ell}, with no undetermined sign.  The
single authority for these conventions is the repository file
\texttt{data/T73\_ENDPOINT\_CONVENTION.json}, which records for each of the
$88$ endpoints its physical identity, owner, orientation, geometric order
(the frozen MWW cut slot), public order (the collar position), THXY index
and pivotal coefficient $\pm q^k$.  The program
\texttt{scripts/build\_t73\_endpoint\_transport.py} derives the monomial
matrix $P(q)$ from it, verifies
$W_{\mathrm{public}}=PW_{\mathrm{geometric}}P^{-1}$ letter by letter along
the actual $45{,}360$-letter cabled word with strand tracking (all eight
orientation patterns occur), and derives $u_0$ and $\ell_0$ from the
canonical coevaluation and the pivotal evaluation instead of reading them;
\texttt{recompute\_t73\_delta3.py} consumes that derivation.  All
omitted factors are invertible $q$-monomials or positive-order
$h$-corrections, and by Lemma~\ref{lem:filtered-cubic} below they do not
affect the divided cubic.

Over $\Q[[h]]$ define
\[
D_h=\ell_h(\rho_h(W)-I)\operatorname{Sh}_h.
\]
This is well defined on the quantum coefficient trace because
\eqref{eq:actualShadow} has already imposed quantum cyclicity.  The exact
Magnus calculation and Darn\'e's Andreadakis theorem give
$W\in\Gamma_3(P_{44})$ \cite[Theorem~6.2]{Darne2021}.  Since every pure
braid generator acts as $I+O(h)$, Lemma~\ref{lem:commutator-order} gives the
one-directional implication
\[
W\in\Gamma_3(P_{44})\ \Longrightarrow\
\rho_h(W)-I\in h^3\operatorname{End}(E_{88});
\]
the converse is not asserted and is not needed.  Independently of the
Andreadakis theorem, direct evaluation of all $7{,}744$ entries of
$\rho_h(W)-I$ verifies the membership on the right.
Thus $D_h$ takes values in $h^3\Q[[h]]$.  Division by $h^3$ is unique, and
reduction modulo $h$ gives a lift-independent ordinary functional
\[
D_3:HH_0(\mathcal C;M_R)\longrightarrow\Q.
\]
Positive-order corrections in the cup, cap, and basis change do not alter
the cubic coefficient (Lemma~\ref{lem:filtered-cubic}).

\begin{lemma}[Filtered cubic pairing]\label{lem:filtered-cubic}
Let $A(h)=\sum_{j\ge0}A_jh^j$ be an endomorphism series of $E_{88}$ with
$A_0=A_1=A_2=0$, and let $u(h)=\sum_k u_kh^k$ and
$\ell(h)=\sum_i\ell_ih^i$ be a vector series and a row series.  Then
\[
[h^3]\,\ell(h)A(h)u(h)=\ell_0A_3u_0,
\]
and the coefficient is unchanged when $(A,u,\ell)$ is replaced by
$(PAP^{-1},Pu,\ell P^{-1})$ for a constant automorphism $P$.
\end{lemma}

\begin{proof}
$[h^3]\ell(h)A(h)u(h)=\sum_{i+j+k=3}\ell_iA_ju_k$, and every term with
$j\le2$ vanishes.  Under simultaneous transport each summand is unchanged.
Both statements are formalized in
\texttt{Smooth4PC/FilteredCubicNaturality.lean}
(\texttt{pairingCoeff\_three\_of\_startsAtThree},
\texttt{cubic\_invariant\_under\_simultaneous\_transport}); the numerical
hypothesis $\rho_h(W)-I\in h^3\operatorname{End}(E_{88})$ is recomputed on
all $7{,}744$ entries by the endpoint transport program, and the cubic
agrees whether it is evaluated with the constant terms, with the full
transported $q$-series, or as $\ell_0A_3u_0$.
\end{proof}

The four mixed-index Burau sign variants
\[
-53824,\quad -53760,\quad -59008,\quad -59072
\]
are nonzero, but they are not the frozen cubic: each mixes two endpoint
coordinate systems.  The frozen public receipt and Lean give
\begin{equation}\label{eq:selectedD3}
[h^3]\,\ell_h(\rho_h(W)-I)u_h=2624\ne0
\end{equation}
in the single collar convention.  Equation~\eqref{eq:selectedD3} is a
statement about the frozen Burau computation; it becomes
$D_3([v_T])=2624$ only when Lemma~\ref{lem:C1} identifies $[v_T]$ with an
actual coefficient-trace class, which is \Open.  Nonvanishing does not
close C or S.

\subsection{The complete two-handle cocone}

\begin{lemma}[All-cable constant term]\label{lem:C2}
For every finite cable level and every sign-preserving cable braid \(b\),
the endpoint operator has constant term equal to its signed permutation.
After standardizing that permutation, the remaining pure-braid operator is
\(I+O(h)\).
\end{lemma}

\begin{proof}
BPW's oriented-cabling action is obtained from the strict BHPW tangle
functor.  At \(q=1\), the braiding on the fundamental representation and its
dual is the ordinary symmetric interchange, including the pivotal
identifications for oppositely oriented factors.  Therefore the constant
term of a braid is exactly the map of its signed endpoint permutation.  A
pure sign-preserving braid has identity permutation and hence constant term
\(I\).  All matrix entries are Laurent polynomials in \(q\), regular at
\(q=1\); after \(q=1+h\), subtracting the constant term leaves a multiple of
\(h\).  The argument is independent of the number of cable copies and so
holds at every finite MWW level.  This statement concerns the endpoint
representation only and does not assert the unknown symmetric-group
factorization of the anti-parallel MWW beta action.
The lemma is applied to the Johnson replacement after Lemma~\ref{lem:C1}.
\end{proof}

Only $D_3$ is asserted to descend after C1.  The remainder of this
subsection is the C2 argument after Lemma~\ref{lem:C1}.

P0c supplies product normals on the reconstruction strands, hence disjoint
parallel levels in the replacement collar.  We first type the construction at every MWW
cable state.  In owner order $(m_2,m_3,r_{xy},r_{yz},r_{zx})$, put
\[
n_y=(42,189,2,2,0),\qquad n_z=(269,1271,2,2,0).
\]
For $r\in\mathbb N^5$, cutting the actual cable $K(r,r)$ along the $y$- and
$z$-cocores gives
\[
p_y(r)=\sum_i r_i n_{y,i},\qquad
p_z(r)=\sum_i r_i n_{z,i}.
\tag{24}
\]
By Lemma~\ref{lem:C1}, once established, parallel copies of the P0 rectangles would give a chain isomorphism $H_r$ of the
same form as (17), with the corresponding endpoint sets.  The
$r_{zx}$ component is not assumed split; its stabilization relations are
the subject of Section~\ref{sec:local-stabilization}.  Since all
copies lie in one fixed tubular product chart, $H_r$ commutes with every
gluing action and with adding a parallel pair.

Let $s_0=e_{m_2}+e_{r_{xy}}$.  If $r\ge s_0$, let $\Omega_r$ be the finite set
of choices of one negative and one positive physical copy of each of $m_2$
and $r_{xy}$.  For $\omega\in\Omega_r$, close every unselected product factor
by the MWW core counit, use the selected 44 passages for the point-push, and
write $D_{r,\omega}$ for the cubic coefficient of the resulting row.  Define
\[
D_r=|\Omega_r|^{-1}\sum_{\omega\in\Omega_r}D_{r,\omega}.
\tag{25}
\]
For a state not above $s_0$, add the missing pairs with the once-dotted MWW
maps and define $D_r$ by pulling (25) back.  Disjoint supports make the order
of these additions immaterial.  Thus $D_r$ is defined on
every summand in
MWW Definition~3.1, including the states below the selected one.

For beta, apply the P0 collar motion simultaneously to all physical copies.
This is the oriented-cabling action of
\cite[Theorem~E]{BeliakovaPutyraWehrli2019}, made strict by BHPW.  Constant
permutations reindex $\Omega_r$.  After a placement is standardized, the
remaining pure braid is \(I+O(h)\) by Lemma~\ref{lem:C2}.  Simultaneous transport of the
operator, cup and cap, together with $\rho(W)-I=O(h^3)$, changes the row only
in order at least four.  Therefore
\[
D_r\beta_i(b)=D_r
\tag{26}
\]
for every $b\in B_{r_i,r_i}$.  This cubic-order argument does not use the
assertion that anti-parallel braid actions factor through a symmetric group,
which MWW leave open
\cite[Remark~3.3]{ManolescuWalkerWedrich2023}.

For a pair addition, split the target placements according to whether the
new pair is selected.  A beta move exchanges a selected new pair with an old
one, so (26) reduces both subsets to the same local core calculation.  On the
whole source that calculation is
\[
(\epsilon\otimes\epsilon)\Delta(1)=0,\qquad
(\epsilon\otimes\epsilon)\Delta(X)=1.
\]
For $r_{zx}$ the same equations are required inside a tubular neighbourhood
of the component, without any split hypothesis; this is the local statement
of Section~\ref{sec:local-stabilization}.  Forgetting the distinguished new pair has constant-size fibers
between placement sets, and the normalizations in (25) cancel their
cardinalities.  The dotted raw degree $+2$ is cancelled by the target cabled
shift $-2$.  Consequently, for every owner and state,
\[
D_{r+e_i}\psi_i^{[0]}=0,\qquad
D_{r+e_i}\psi_i^{[1]}=D_r.
\tag{27}
\]
Below $s_0$ the second equality is the definition; the Frobenius identities
and disjoint-support interchange prove the first equality and the commutation
of all owner squares.  The algebra of this statewise construction is
formalized in \texttt{Smooth4PC/ReynoldsCableCocone.lean}: for every state
$r\in\mathbb N^n$, the Reynolds average is invariant under a braid action
that permutes placements, the two pair-addition relations follow from
placement-wise relations with constant fibres, and the extension of $D_r$
to the states below $s_0$ along iterated once-dotted additions is
independent of the common upper bound and satisfies the $\psi^{[1]}$
relation at every state.  Every geometric input of that file (that the
actual MWW beta action permutes placements, that the actual pair additions
forget a distinguished pair with constant fibres and commute, and the
descent above $s_0$) is a hypothesis there, and for the actual cabled
source these hypotheses are part of the \Open\ Hypothesis~\ref{hyp:P1}.

Once (26)--(27) are established on the actual cabled source, together with
the coefficient trace relation of (17), they prove on the whole cabled
direct sum
\[
D_3(\beta_i(b)x-x)=0,\quad
D_3(\psi_i^{[0]}x)=0,\quad
D_3(\psi_i^{[1]}x-x)=0,
\]
and the quotient universal property with MWW Theorem~3.2 then gives
$\overline D_3:\mathcal S^2_0(W_2;\Q)\to\Q$.  At present (17) and
(26)--(27) are verified only on the control collar, so this descent is part
of the \Open\ Hypothesis~\ref{hyp:P1}.

\subsection{Local stabilization without a split hypothesis}
\label{sec:local-stabilization}

For a framed component $K_i$ and a cable state $r$, let $\chi_{i,r}$ denote
the local closure of the $r_i$ parallel copies of $K_i$ inside a tubular
neighbourhood of $K_i$, defined in the pre-quotient foam category by the
undotted and once-dotted sphere evaluations.  The statement needed for the
$\psi$ relations of every owner, including $r_{zx}$, is
\[
\chi_{i,r+e_i}\circ\psi_i^{[0]}=0,\qquad
\chi_{i,r+e_i}\circ\psi_i^{[1]}=\chi_{i,r},
\]
with no assumption that $K_i$ is split from the rest of the link.  This
local stabilization--retraction statement is \Open.  The earlier text
inferred a split unknot from the empty free word of the transported
$r_{zx}$; that inference is withdrawn, since a free-group word does not
determine an embedded component.

\subsection{C conclusion and the degree-$494$ square}

The MWW one- and two-handle maps fit into the commuting diagram
\[
\begin{array}{ccccccc}
M_R(T_1,T_1)&\xrightarrow{H_{T_1,T_1}}&
\operatorname{End}(U)\{-44\}\otimes A^{\otimes227}
&\xrightarrow{\operatorname{id}\otimes\epsilon^{\otimes227}}&
\operatorname{End}(U)&\xrightarrow{d_3}&\Q\\
\big\downarrow{\mathrm{Glue}_{1h}}&&&&&&\big\Vert\\
\mathcal S^2_0(W_1;K(s_0,s_0))\{315\}
&\xrightarrow{\ \iota_{s_0}\{-4\}\ }&
\mathcal S^{2,0}_0(W_1;K)/(\beta,\psi)
&\xrightarrow[\cong]{\ \Phi\ }&\mathcal S^2_0(W_2)
&\xrightarrow{\overline D_3}&\Q.
\end{array}
\tag{28}
\]
The left vertical map and its shift are MWW Theorem~4.7, and the bottom
isomorphism is MWW Theorem~3.2.  The upper selected element is
$\operatorname{Id}_U\otimes X^{\otimes227}$.  Its successive degrees are
\[
-44+227=183,\qquad183+315=498,\qquad498-4=494.
\]
Under Hypothesis~\ref{hyp:P1}, diagram (28) sends it to a class $x_2$
with $\overline D_3(x_2)=2624\ne0$ on the Johnson replacement collar.

\begin{proposition}[Status of C]\label{thm:Cdischarge}
Hypothesis~\ref{hyp:P1} is \Open\ for the Johnson replacement.
Established: the representable $HH_0$ reduction, Lemma~\ref{lem:C2} for the
endpoint representation, Lemma~\ref{lem:filtered-cubic}, and the frozen
Burau value \eqref{eq:selectedD3}.  Not established: the actual
product-ribbon isotopy and chain-level action squares of
Lemma~\ref{lem:C1}; the statewise definition of $D_r$ by placement averages
and its relations (26)--(27) on every cable state of the actual source; the
local stabilization statement of Section~\ref{sec:local-stabilization}; and
hence the descent of $D_3$ through the complete two-handle cocone and the
nonvanishing of the selected degree-$494$ class after all MWW two-handle
relations.
\end{proposition}

\begin{proof}
The established items are the replayed certificates and the Lean files
named in Appendix~\ref{sec:reproducibility-extended}.  The remaining items
require the actual cut tangle, which is not part of the current
certificates.
\end{proof}

\paragraph{C/S firewall.}
The $B$-external naturality obtained in C concerns boundary cobordisms in
the coefficient comparison.  It does not identify the local action of an
essential sphere in $\partial W_2$ with an ordinary dotted-sphere evaluation;
that is the separate calculation in Lemma~\ref{lem:Sendpoint}.
